\section{Future work}\label{sec:future_work}

This section outlines potential future work and improvements that can be made to the proposed system.
While the current implementation has demonstrated its effectiveness in crawling, processing, and analyzing e-commerce data, there are several areas where further enhancements can be made.

\subsection{Multi-layer crawling architecture}\label{subsec:multi-layer-crawling-architecture}
The current crawler architecture relies on a single queue to manage the crawling tasks.
While AWS SQS' queues do allow for redrive policies to handle failed messages, the retries are currently handled by the same crawler that failed to process it in the first place.
Upon three retries, the message is sent to a dead-letter queue, and the failed page is not crawled again until the next crawl cycle.
This can be improved by implementing a multi-layer crawling architecture, where failed pages are sent to a separate queue that is processed by a different set of crawlers, which can be configured to handle specific types of failures, such as network errors, timeouts, or server errors.
They could even be crawled by a crawler instance running on a separate AWS region, which would allow for a different network path to be used, and could potentially avoid the issues that caused the failure in the first place.
This would allow for more efficient handling of failed pages, and would increase the overall success rate of the crawling process.

\subsection{Dynamic crawling frequency}\label{subsec:dynamic-crawling-frequency}
The current implementation of the system spawns a new crawler instance every 3 minutes, which crawls 100 pages.
While this is a reasonable crawling frequency for the current set of pages, if the dataset were to grow (or shrink) significantly, this crawling frequency may not be optimal.
This can be improved by implementing a dynamic crawling frequency, where the crawling frequency is adjusted based number of pages in the queue, and the number of pages to be crawled by each instance is adjusted based on the average time taken to crawl a page.

\subsection{Memory optimization on classification and extraction inference}\label{subsec:memory-optimization-on-classification-and-extraction-inference}
The current implementation of the jobs used to generate datasets and perform inference for the page classification and information extraction models are run on AWS Sagemaker instances which require a large amount of memory.
If the dataset were to grow significantly, this could lead to memory issues, and the jobs may fail to complete successfully.
This can be improved by implementing memory optimization techniques, such as using smaller batch sizes, or using model quantization techniques to reduce the memory footprint of the models, which would allow for more efficient use of memory and reduce the risk of memory-related issues during inference.

\subsection{Cost optimization}\label{subsec:cost-optimization}
While the current implementation of the system is cost-effective, there are several areas where further cost optimization can be achieved.
\begin{itemize}
    \item Using AWS reserved instances and saving plans for the crawler instances and the RDS database, which can provide significant cost savings compared to on-demand instances.
    \item Minimizing the frequency of page classification.
    Currently, all pages are classified once per crawl cycle.
    However, it is very unlikely that the classification of a page will change between crawl cycles, and therefore, it may be more cost-effective to only classify pages that have changed since the last crawl cycle, or to classify at a more infrequent interval.
    \item Dynamic frequency of information extraction.
    For products that are marked as out of stock, it is very unlikely that the product information will change between crawl cycles, and therefore, it may be more cost-effective to only extract information for products that are in stock, or to extract information at a more infrequent interval.
    Even for pages that have had static pricing for a long period of time, it would be unlikely that the price to have changed in a single crawl cycle.
    It may be effective to reduce the frequency of information extraction for products that have had static pricing for a long period of time, and only extract information for products that have had price changes recently.
\end{itemize}

\subsection{Study on pricing practices by online retailers}\label{subsec:study-on-pricing-practices-by-online-retailers}
The dataset generated by this project can be repurposed to conduct a study on pricing practices by online retailers.
This can be done by analyzing the price history of products across different e-commerce platforms, and identifying patterns and trends in pricing practices, such as price discrimination, dynamic pricing, and fake discounts.
This can also be used to identify the impact of external factors, such as holidays, sales events, and market trends, on pricing practices by online retailers, which will be useful in predicting future pricing trends and identifying potential opportunities for consumers to make informed purchasing decisions.